% Stile grafico: colori, intestazioni, titoli, ambienti e comandi del documento.

% --- Colori ---------------------------------------------------------------
\definecolor{unina}{RGB}{0,62,126}
\definecolor{uninachiaro}{RGB}{226,234,244}
\definecolor{grigiotesto}{RGB}{90,90,90}
\definecolor{rossotodo}{RGB}{176,28,28}

% --- Collegamenti ---------------------------------------------------------
\hypersetup{
  colorlinks=true,
  linkcolor=unina,
  urlcolor=unina,
  citecolor=unina,
  pdftitle={\progetto{} -- \titolodoc},
  pdfauthor={\autoreuno},
  pdfsubject={\insegnamento},
}

% --- Titoli dei capitoli e delle sezioni ----------------------------------
\titleformat{\chapter}[display]
  {\normalfont\color{unina}}
  {\large\scshape Capitolo \thechapter}
  {4pt}
  {\Huge\bfseries}
  [\vspace{4pt}{\color{unina}\titlerule[1pt]}]
\titlespacing*{\chapter}{0pt}{0pt}{28pt}

\titleformat{\section}{\normalfont\Large\bfseries\color{unina}}{\thesection}{0.8em}{}
\titleformat{\subsection}{\normalfont\large\bfseries}{\thesubsection}{0.8em}{}

% --- Intestazioni e piè di pagina -----------------------------------------
\pagestyle{fancy}
\fancyhf{}
\renewcommand{\chaptermark}[1]{\markboth{#1}{}}
\fancyhead[L]{\small\color{grigiotesto}\leftmark}
\fancyhead[R]{\small\color{grigiotesto}\progetto{} -- \titolodoc}
\fancyfoot[L]{\small\color{grigiotesto}Versione \versionedoc{} -- \statodoc}
\fancyfoot[R]{\small\thepage}
\renewcommand{\headrulewidth}{0.4pt}
\renewcommand{\footrulewidth}{0.4pt}
\fancypagestyle{plain}{%
  \fancyhf{}
  \fancyfoot[R]{\small\thepage}
  \renewcommand{\headrulewidth}{0pt}
  \renewcommand{\footrulewidth}{0pt}
}

% --- Liste e tabelle ------------------------------------------------------
\setlist[itemize]{topsep=3pt, itemsep=2pt}
\setlist[enumerate]{topsep=3pt, itemsep=2pt}
\captionsetup{font=small, labelfont={bf,color=unina}}
\renewcommand{\arraystretch}{1.2}
\newcolumntype{L}{>{\raggedright\arraybackslash}X}

% --- Comandi --------------------------------------------------------------
% Parte ancora da scrivere: compare in rosso nel PDF.
\newcommand{\todo}[1]{\textcolor{rossotodo}{\textbf{[DA COMPLETARE: #1]}}}

% Stile dei blocchi nei diagrammi TikZ.
\tikzset{
  box/.style={draw=unina, thick, rounded corners, fill=uninachiaro,
              minimum width=2.4cm, minimum height=1cm, align=center, font=\small},
  freccia/.style={-{Latex}, thick, color=unina!80!black},
}

% Diagrammi delle architetture: componenti, oggetti di Kubernetes, messaggi.
\definecolor{accento}{RGB}{204,102,0}
\tikzset{
  pod/.style={draw=unina, thick, rounded corners=2pt, fill=white,
              minimum width=1.7cm, minimum height=0.8cm, align=center, font=\small},
  copie/.style={pod, double copy shadow={shadow xshift=3pt, shadow yshift=3pt,
                draw=unina, fill=white}},
  kube/.style={draw=unina, fill=unina, text=white, rounded corners=2pt,
               minimum height=0.7cm, inner sep=3pt, align=center, font=\scriptsize\bfseries},
  db/.style={cylinder, shape border rotate=90, aspect=0.22, draw=unina, thick,
             fill=uninachiaro, minimum width=2cm, minimum height=1.3cm,
             align=center, font=\small},
  coda/.style={rectangle split, rectangle split horizontal, rectangle split parts=4,
               rectangle split empty part width=1.4mm, draw=unina, thick,
               fill=uninachiaro, minimum height=0.75cm},
  toxi/.style={draw=grigiotesto, densely dashed, rounded corners=2pt, fill=white,
               minimum height=0.7cm, inner sep=3pt, font=\scriptsize, text=grigiotesto},
  guastato/.style={draw=accento, text=accento, fill=accento!8},
  gruppo/.style={draw=unina!55, dashed, rounded corners=3pt, inner sep=7pt},
  cluster/.style={draw=grigiotesto!60, rounded corners=4pt, inner sep=9pt, fill=black!2},
  etichetta/.style={font=\scriptsize, text=grigiotesto, align=center},
  risposta/.style={-{Latex}, thick, dashed, color=unina!80!black},
  sonda/.style={-{Latex}, densely dotted, thick, color=grigiotesto},
  vita/.style={draw=grigiotesto!70, dashed},
}

% --- Decisioni architetturali (ADR) ---------------------------------------
% \begin{decisione}{ID}{Titolo}{Stato}{Data} ... \end{decisione}
% Ogni decisione è una sezione del capitolo "Decisioni architetturali".
\newenvironment{decisione}[4]{%
  \section{#1 -- #2}
  \noindent
  \begin{tcolorbox}[colback=uninachiaro, colframe=unina, boxrule=0.6pt, arc=2pt,
                    left=6pt, right=6pt, top=4pt, bottom=4pt]
    \small
    \begin{tabularx}{\linewidth}{@{}lL@{}}
      \textbf{Identificativo} & #1 \\
      \textbf{Stato}         & #3 \\
      \textbf{Data}          & #4 \\
    \end{tabularx}
  \end{tcolorbox}
}{}

% Sottotitoli fissi dentro una decisione.
\newcommand{\contesto}{\paragraph{Contesto.}}
\newcommand{\scelta}{\paragraph{Decisione.}}
\newcommand{\conseguenze}{\paragraph{Conseguenze.}}
